\RequirePackage{pdfmanagement}
\SetKeys[document/metadata]{pdfversion=1.7, pdfstandard=A-2u, lang=en-GB}
\input{preamble/pdfA-compliance/glyphtounicode}
\documentclass[11pt,a4paper]{report}

\newcommand*{\PaperTitle}{Quadratic \((n,n+2,n)\) mass-action networks with \(n+1\) positive nondegenerate equilibria}
\newcommand*{\PaperTitlePlain}{Quadratic (n,n+2,n) mass-action networks with n+1 positive nondegenerate equilibria}
\newcommand*{\PaperShortTitle}{Quadratic \((n,n+2,n)\) networks with \(n+1\) equilibria}
\newcommand*{\PaperSubtitle}{chemistry-open-problems \textperiodcentered{} 18 October 2026 \textperiodcentered{} preliminary}
\newcommand*{\PaperCredits}{Solver: Claude Opus 5.5 (effort xhigh), orchestrating agent \textperiodcentered{} Independent verification: pending}
\newcommand*{\PaperKeywords}{chemical reaction networks, multistationarity, mass-action kinetics, Gale duality, quadratic networks}
\input{chemop-preamble}
\addbibresource{references.bib}

\newcommand{\Net}{\mathcal N}
\providecommand{\im}{\operatorname{im}}
\providecommand{\spn}{\operatorname{span}}

\begin{document}
\makepapertitle

\begin{paperabstract}
Banaji and Feliu proved that a full-rank quadratic mass-action network with \(n\) species and \(n+2\) distinct
sources has at most \(n+1\) positive nondegenerate equilibria, and wrote that they could not find a quadratic
\((3,5,3)\) network with four, leaving open whether the bound is attained for \(n\ge3\). We give, for every \(n\ge1\), an
explicit quadratic \((n,n+2,n)\) network whose sources are \(0,2X_1,X_1+X_2,\dots,X_{n-1}+X_n,X_n\) and which has exactly
\(n+1\) positive nondegenerate equilibria for suitable rate constants. Its kernel can be prescribed freely, which reduces
the count to the zeros of an explicit one-variable function; geometrically separated kernel ratios then force the
maximal number of equilibria. For \(n=1\) the network is \(0\to X,\ X\to0,\ 2X\to3X\), which shows that the bound is
attained for \(n=1\) as well, contrary to Remark~5.6 of the same paper.
\end{paperabstract}

\section*{The question}

A network is \emph{quadratic} if all its sources have molecularity at most two (products are unrestricted), and an
\((n,m,r)\) network has \(n\) species, \(m\) reactions and rank \(r\) \autocite{BanajiFeliu2026}. Theorem~5.5 of
\autocite{BanajiFeliu2026} states that a dynamically nontrivial quadratic \((n,n+k,n)\) network with \(n+2\) distinct
sources which is not source deficient admits no more than \(n+1\) positive nondegenerate equilibria. In Remark~5.8 the
authors write (arXiv:2409.06877v5, p.~41):

\begin{openquote}
``However, we were not able to find a quadratic \((3,5,3)\) network admitting four positive nondegenerate equilibria,
and thus achieving the bound in Theorem 5.5. We leave open the question of whether the bound in Theorem 5.5 can be
achieved by quadratic \((n, n+2, n)\) networks for \(n\ge3\).''
\end{openquote}

\begin{theorem}\label{thm:main}
For every \(n\ge1\) there is a quadratic \((n,n+2,n)\) mass-action network with \(n+2\) distinct sources, dynamically
nontrivial and not source deficient, which for suitable rate constants has exactly \(n+1\) positive equilibria, all of
them nondegenerate. Hence the bound of \autocite[Theorem~5.5]{BanajiFeliu2026} is attained for every \(n\ge1\) (with
\(k=2\)).
\end{theorem}

\begin{remark}[note=on Remark~5.6 of \autocite{BanajiFeliu2026}]
Remark~5.6 of \autocite{BanajiFeliu2026} states that for \(n=1\) the bound is not optimal, because ``by [5, Lemma 20],
quadratic rank-one networks do not admit multiple positive nondegenerate equilibria''. Lemma~20 of the cited paper
\autocite{BBH2024ND} says the opposite: a quadratic rank-one network admits multiple positive nondegenerate equilibria
if and only if it contains one of four small networks, the first being \(0\to X,\ X\to0,\ 2X\to3X\). The statement
quoted in Remark~5.6 holds for \emph{bimolecular} rank-one networks \autocite[Corollary~19]{BBH2024ND}. Indeed
\(\dot x=\kappa_1-\kappa_3x+\kappa_2x^2\) has the two positive simple roots \((3\pm\sqrt5)/2\) for
\(\kappa=(1,1,3)\). This network is \(\Net_1\) below.
\end{remark}

\section{The networks}

Fix \(n\ge1\), species \(X_1,\dots,X_n\), and positive integers \(\alpha_2,\dots,\alpha_{n+1}\) and
\(\beta_2,\dots,\beta_{n+1}\). Write \(\alpha=(1,0,\alpha_2,\dots,\alpha_{n+1})\) and \(\beta=(0,1,\beta_2,\dots,\beta_{n+1})\).
The network \(\Net_n(\alpha,\beta)\) has the \(n+2\) reactions
\[
\begin{aligned}
R_0&:\ 0\ \longrightarrow\ \textstyle\sum_{i=1}^n\alpha_{i+1}X_i, &
R_1&:\ 2X_1\ \longrightarrow\ 2X_1+\textstyle\sum_{i=1}^n\beta_{i+1}X_i,\\
R_{i+1}&:\ X_i+X_{i+1}\ \longrightarrow\ X_{i+1}\quad(1\le i\le n-1), &
R_{n+1}&:\ X_n\ \longrightarrow\ 0 .
\end{aligned}
\]
Its sources \(\sigma_0=0\), \(\sigma_1=2X_1\), \(\sigma_{i+1}=X_i+X_{i+1}\), \(\sigma_{n+1}=X_n\) are distinct and quadratic,
and the products are nonnegative. Let \(\Gamma\in\mathbb Z^{n\times(n+2)}\) be the stoichiometric matrix (columns indexed
\(0,\dots,n+1\)) and \(S\in\mathbb Z^{(n+2)\times n}\) the source matrix. The \(i\)-th row of \(\Gamma\) is
\(\alpha_{i+1}e_0+\beta_{i+1}e_1-e_{i+1}\).

\begin{lemma}\label{lem:structure}
\(\Net_n(\alpha,\beta)\) is a full-rank \((n,n+2,n)\) network with \(\ker\Gamma=\spn(\alpha,\beta)\) and
\(\ker\Gamma\cap\mathbb R^{n+2}_{>0}=\{s\alpha+t\beta:s,t>0\}\); in particular it is dynamically nontrivial. Moreover
\(\operatorname{rank}[S\,|\,\mathbf 1]=n+1\), and the left kernel of \([S\,|\,\mathbf 1]\) is spanned by
\(w=(w_0,w_1,\dots,w_{n+1})\) with
\[
w_1=1,\qquad w_j=2(-1)^{j-1}\ (2\le j\le n+1),\qquad w_0=-\textstyle\sum_{j\ge1}w_j .
\]
\end{lemma}

\begin{proof}
Row \(i\) applied to \(\alpha\) gives \(\alpha_{i+1}\cdot1+\beta_{i+1}\cdot0-\alpha_{i+1}=0\), and to \(\beta\) gives
\(0+\beta_{i+1}-\beta_{i+1}=0\). The entries \(-1\) sit in distinct columns \(2,\dots,n+1\), so \(\operatorname{rank}\Gamma=n\) and
\(\ker\Gamma\) is the two-dimensional space \(\spn(\alpha,\beta)\). Since \(\alpha_1=\beta_0=0\) and all other entries are
positive, \(s\alpha+t\beta>0\) iff \(s,t>0\). The sources \(X_n,X_{n-1}+X_n,\dots,X_1+X_2\) span \(\mathbb R^n\) and \(\sigma_0=0\),
so \(\operatorname{rank}[S|\mathbf1]=n+1\). Finally \(\sum_jw_j\sigma_j=0\): the coefficient of \(X_1\) is \(2w_1+w_2=2-2=0\), and
that of \(X_i\) (\(2\le i\le n\)) is \(w_i+w_{i+1}=0\); and \(\sum_jw_j=0\) by the choice of \(w_0\).
\end{proof}

\section{Equilibria as roots of one function}

Put \(v(u)=\alpha+u\beta\) for \(u>0\), so \(v_0=1\), \(v_1=u\), \(v_j=\alpha_j+u\beta_j\) (\(j\ge2\)), and
\[
f(u)=\sum_{j=0}^{n+1}w_j\log v_j(u)=\log u+\sum_{j=2}^{n+1}2(-1)^{j-1}\log(\alpha_j+\beta_ju).
\]

\begin{lemma}\label{lem:bijection}
Let \(\kappa\in\mathbb R^{n+2}_{>0}\) and \(c=\sum_jw_j\log\kappa_j\). The positive equilibria of \(\Net_n(\alpha,\beta)\) with rate
constants \(\kappa\) are in bijection with the solutions \(u>0\) of \(f(u)=c\).
\end{lemma}

\begin{proof}
\(x>0\) is an equilibrium iff the rate vector \(\rho_j=\kappa_jx^{\sigma_j}\) lies in \(\ker\Gamma\cap\mathbb R^{n+2}_{>0}\), i.e.\
\(\rho=s\,v(u)\) for unique \(s,u>0\) (\zcref{lem:structure}). Taking logarithms, this is the linear system
\(\sigma_j\cdot\log x-\log s=\log v_j(u)-\log\kappa_j\) (\(0\le j\le n+1\)) in the \(n+1\) unknowns \((\log x,\log s)\), whose
coefficient matrix \([S|-\mathbf1]\) has rank \(n+1\) and left kernel \(\spn(w)\). It is solvable iff
\(\sum_jw_j(\log v_j(u)-\log\kappa_j)=0\), i.e.\ \(f(u)=c\), and then its solution is unique.
\end{proof}

\begin{lemma}\label{lem:nondeg}
The equilibrium corresponding to \(u\) is nondegenerate (its Jacobian is nonsingular) iff \(f'(u)\neq0\).
\end{lemma}

\begin{proof}
The Jacobian is \(J=\Gamma D_\rho SD_{1/x}\) with \(\rho=s\,v(u)\). As \(S\) is injective, \(J\) is singular iff
\(\im S\cap D_\rho^{-1}\ker\Gamma\neq\{0\}\). The left kernel of \(S\) is two-dimensional and contains \(w\); since
\(\mathbf1\notin\im S\), it also contains some \(w'\) with \(w'\cdot\mathbf1=1\), and \(\im S=\{q:w\cdot q=w'\cdot q=0\}\). For
\(a\alpha+b\beta=a\,v(u)+(b-au)\beta\) we get
\(q=D_\rho^{-1}(a\alpha+b\beta)=s^{-1}\bigl(a\mathbf1+(b-au)(\beta_j/v_j)_j\bigr)\), hence
\(w\cdot q=s^{-1}(b-au)f'(u)\) and \(w'\cdot q=s^{-1}\bigl(a+(b-au)\textstyle\sum_jw'_j\beta_j/v_j\bigr)\). If \(f'(u)\ne0\),
both vanish only for \(a=b=0\); if \(f'(u)=0\), \(b-au=1\), \(a=-\sum_jw'_j\beta_j/v_j\) gives a nonzero \(q\in\im S\).
\end{proof}

\begin{lemma}\label{lem:degree}
For any \(c\), \(f-c\) has at most \(n+1\) zeros in \((0,\infty)\), and if it has \(n+1\) they are all simple.
\end{lemma}

\begin{proof}
\(f'(u)=1/u+\sum_{j=2}^{n+1}w_j\beta_j/(\alpha_j+\beta_ju)\). Multiplying by \(u\prod_{j\ge2}(\alpha_j+\beta_ju)\) gives a
polynomial of degree at most \(n\) whose coefficient of \(u^n\) is \(\prod_j\beta_j\cdot(1+\sum_{j\ge2}w_j)=\pm\prod_j\beta_j\neq0\).
So \(f'\) has at most \(n\) zeros. By Rolle's theorem, \(n+1\) zeros of \(f-c\) force \(n\) zeros of \(f'\) strictly between
them, so none of the zeros of \(f-c\) is a zero of \(f'\).
\end{proof}

\section{Geometrically separated kernels}

Fix \(R\ge64\) and take \(\alpha_j=R^{j-1}\), \(\beta_j=1\) for \(2\le j\le n+1\); write \(L=\log R\). Then
\(f(u)=\log u+\sum_{m=1}^{n}2(-1)^m\log(R^m+u)\).

\begin{proposition}\label{prop:signs}
Let \(M_1=f\)'s ``main value'' at \(u=R\), namely \(M_1=L\bigl(1+2\sum_{m=1}^n(-1)^m\max(m,1)\bigr)\), and \(c=M_1-L/2\). Then
\(f(R^i)-c>0\) for odd \(i\) and \(f(R^i)-c<0\) for even \(i\) (\(1\le i\le n\)). Moreover \(f(u)\to-\infty\) as \(u\to0^+\), and
\(f(u)\to-\infty\) (\(n\) odd) or \(+\infty\) (\(n\) even) as \(u\to\infty\). Consequently \(f-c\) has exactly \(n+1\) zeros in
\((0,\infty)\), all simple.
\end{proposition}

\begin{proof}
For \(1\le i\le n\), \(\log(R^m+R^i)=\max(m,i)L+\log(1+R^{-|m-i|})\). Hence
\(f(R^i)=M_i+2(-1)^i\log2+\varepsilon_i\) with \(M_i=L\bigl(i+2\sum_{m=1}^n(-1)^m\max(m,i)\bigr)\) and
\(|\varepsilon_i|\le\sum_{d\ge1}2\cdot2\log(1+R^{-d})\le 4/(R-1)\) (each \(d=|m-i|\ge1\) occurs at most twice). A direct
computation gives \(M_{i+1}-M_i=L\bigl(1+2\sum_{m=1}^i(-1)^m\bigr)\), which is \(-L\) for odd \(i\) and \(+L\) for even \(i\); so
\(M_i=M_1\) for odd \(i\) and \(M_i=M_1-L\) for even \(i\). Therefore
\[
f(R^i)-c=\begin{cases}L/2-2\log2+\varepsilon_i, & i\text{ odd},\\ -(L/2-2\log2)+\varepsilon_i, & i\text{ even},\end{cases}
\]
and \(L/2-2\log2\ge\log2>4/63\ge|\varepsilon_i|\) for \(R\ge64\). As \(u\to0^+\), \(f(u)=\log u+O(1)\to-\infty\). As
\(u\to\infty\), \(f(u)=(1+\sum_{m=1}^n2(-1)^m)\log u+o(\log u)\) with \(1+\sum_{m=1}^n 2(-1)^m=-1\) (\(n\) odd) or \(+1\) (\(n\) even).
So \(f-c\) takes alternating signs at \(n+2\) points \(0^+<R<R^2<\dots<R^n<\infty\), has at least \(n+1\) zeros, hence exactly
\(n+1\), all simple, by \zcref{lem:degree}.
\end{proof}

\begin{proof}[of \zcref{thm:main}]
Take \(\Net_n(\alpha,\beta)\) with \(\alpha_j=64^{j-1}\), \(\beta_j=1\), and rate constants \(\kappa_j=1\) for \(j\ne1\) and
\(\kappa_1=e^c\) with \(c\) as in \zcref{prop:signs} (\(w_1=1\), so \(\sum_jw_j\log\kappa_j=c\)). By
\zcref{lem:bijection,prop:signs,lem:nondeg} the network has exactly \(n+1\) positive equilibria, all nondegenerate. By
\zcref{lem:structure} it is a full-rank quadratic \((n,n+2,n)\) network with \(n+2\) distinct sources, dynamically nontrivial
and not source deficient. The set of admissible \(c\) is open, so \(\kappa_1\) may be taken rational.
\end{proof}

\section{Small explicit examples}

The separation \(R=64\) is convenient for a uniform proof but far from necessary. For \(n=3\), the network
\(\Net_3(\alpha,\beta)\) with \((\alpha_2,\alpha_3,\alpha_4)=(1,3,16)\) and \((\beta_2,\beta_3,\beta_4)=(14,3,1)\),
\[
0\to X_1+3X_2+16X_3,\quad 2X_1\to16X_1+3X_2+X_3,\quad X_1+X_2\to X_2,\quad X_2+X_3\to X_3,\quad X_3\to0,
\]
whose products have molecularity at most 20, has four positive nondegenerate equilibria for a suitable rational
\(\kappa_1\) (all other rate constants equal to 1); the certificate is described below. A quadratic \((3,5,3)\) network
outside this family, found by a computer search, is \(Z\to78Y,\ 2Z\to Z,\ 2Y\to X+Y,\ X\to0,\ X+Y\to24X+20Z\); it also has
four positive nondegenerate equilibria (\texttt{scripts/other\_example\_n3.py}).

\section{Computer verification}

The proof above needs no computer. As a check, \texttt{scripts/certify\_family.py} builds \(\Net_n(\alpha,\beta)\) for the
small \(n=3\) example, for \(0\to X,\ 2X\to3X,\ X\to0\), and for the \(R=64\) family with \(n=1,\dots,8\), and verifies in exact
arithmetic: quadratic distinct sources, nonnegative products, \(\operatorname{rank}\Gamma=n\),
\(\operatorname{rank}[S|\mathbf1]=n+1\), \(\Gamma\alpha=\Gamma\beta=0\), the dependency \(w\), that \(f'\) has exactly \(n\) positive
zeros, and that for an explicit rational \(C\) the polynomial
\(\prod_{w_j>0}v_j^{w_j}-C\prod_{w_j<0}v_j^{-w_j}\) has exactly \(n+1\) positive roots, all simple (Sturm sequences and
\(\gcd(P,P')=1\)); it also computes the corresponding equilibria to 80 digits and checks \(\Gamma\rho(x)=0\) and
\(\det J(x)\neq0\) (101 checks, all pass).

\paperreferences
\end{document}
